\uuid{BQFK}
\titre{Équation droite/plan}
\auteur{bodin ; hanani ; mzari}
\organisation{exo7}
\licence{CC BY-NC-SA 4.0}
\datecreate{2020-01-20}
\niveau{L1}
\module{Géométrie}
\chapitre{Géométrie du plan}
\sousChapitre{Distance, norme, produit scalaire}
\difficulte{3}
\theme{équation droite, plan}
\exo7qcmid{lille-1.geometrie-plane.028}

\contenu{
\begin{question}
\qclassification{140.01,140.02}
\qtags{motcle=équation droite/plan}
Dans le plan rapporté à un repère orthonormé, on considère la droite $\Delta$ d'équation $x=y$ et on note $D$ une droite perpendiculaire à $\Delta$ et qui est à distance $0$ de l'origine. Une représentation paramétrique de $D$ est
\begin{answers}
    \bad{$D\; :\; x=t,\; y=t$ et $t\in \Rr$}
    \good{$D\; :\; x=t,\; y=-t$ et $t\in \Rr$}
    \bad{$D\; :\; x=3t,\; y=3t$ et $t\in \Rr$}
    \good{$D\; :\; x=-2t,\; y=2t$ et $t\in \Rr$}
\end{answers}
\begin{explanations}
Le vecteur $\vec{n}=(1,-1)$ est normal à $\Delta$, il dirige $D$. Or $\mbox{d}(O,D)=0\Rightarrow O\in D$. Donc $D$ est la droite passant par $O$ est dirigée par $\vec{n}$.
\end{explanations}
\end{question}
}
